\documentclass[10pt]{article}

\usepackage[margin=0.85in]{geometry}
\usepackage[T1]{fontenc}
\usepackage{lmodern}
\usepackage{microtype}
\usepackage{booktabs}
\usepackage{tabularx}
\usepackage{array}
\usepackage{graphicx}
\usepackage{xcolor}
\usepackage[numbers,sort&compress]{natbib}
\usepackage[hidelinks]{hyperref}

\definecolor{circteal}{HTML}{176B87}
\newcolumntype{Y}{>{\raggedright\arraybackslash}X}
\newcommand{\placeholderfigure}[1]{%
  \fbox{\parbox[c][1.55in][c]{0.94\linewidth}{\centering\small\color{gray}%
  Original image unavailable in supplied manuscript text.\\[0.4em]#1}}}

\title{\textbf{Coordination Before Collision: Protocol-Level Governance for Multi-Agent Clinical AI}}
\author{%
Bruce Changlong Xu$^{1}$, Ayush Chopra$^{2,3}$, Alexander J. Ryu$^{4}$\\[0.5em]
\small $^{1}$Stanford University, Palo Alto, California, USA\\
\small $^{2}$Massachusetts Institute of Technology, Cambridge, Massachusetts, USA\\
\small $^{3}$Project Iceberg at MIT, Cambridge, Massachusetts, USA\\
\small $^{4}$Mayo Clinic, Rochester, Minnesota, USA\\[0.5em]
\small Correspondence: Alexander J. Ryu}
\date{}

\begin{document}
\maketitle

\begin{abstract}
Autonomous AI agents are beginning to enter clinical workflows, helping to schedule appointments,
manage medication regimens, and route referrals. Yet healthcare interoperability infrastructure was
designed primarily for episodic, human-initiated data exchange rather than continuous coordination
among autonomous actors. This architectural mismatch creates interaction risks: harms that may emerge
when agents optimize independently across shared patients, workflows, and clinical resources. Because
this class of risk arises from agent interaction rather than individual agent behavior alone, it cannot
be fully assessed through single-agent evaluation and may also require population-level evaluation.
Existing healthcare standards, including FHIR, SMART on FHIR, and CDS Hooks, and general-purpose agent
protocols, including A2A, MCP, and ACP, provide foundations for data exchange, authorization, tool
access, and agent interoperability. Clinical multi-agent deployment may additionally require controls
for identity, attribution, scope enforcement, dependency tracking, escalation, and population-level
coordination. We introduce CIRC, Clinical Infrastructure for Reliable Coordination, a conceptual
governance framework intended to stimulate protocol design rather than establish a standard. CIRC
describes four coordination levels, from isolated agents operating through existing EHR interfaces to
system-aware populations coordinating over aggregate resource and acuity signals, and identifies five
protocol primitives associated with those levels. Using a hypothetical oncology vignette, we illustrate
how interaction risks could emerge at each level and how richer infrastructure could create decision
points for escalation, review, and failure containment. We conclude with implications for standards
bodies, deploying institutions, and regulators, including simulation-based stress testing before live
deployment and post-deployment monitoring in clinical practice.
\end{abstract}

\section{Introduction}

Clinical artificial intelligence (AI) is shifting from decision support toward the possibility of
autonomous action by AI agents. Traditional medical AI systems advise clinicians, who remain
responsible for clinical decisions and downstream actions. A new generation of agents may schedule
appointments, initiate referrals, manage prior authorization, and coordinate multispecialty care with
minimal human intervention.\cite{moritz2025coordinated,borkowski2025multiagent,taylor2025agents} This
transition from advisory to agentic systems is a qualitative change in how AI participates in
healthcare delivery, one that current interoperability infrastructure was not designed to support.

The challenge is architectural mismatch. Healthcare's dominant standards, including FHIR, SMART on
FHIR, and CDS Hooks, were designed primarily for episodic data exchange between human-operated
systems, with the assumption that humans initiate or mediate clinical actions.\cite{hl7fhir2019,
hl7cdshooks2018,mandel2016smart} Autonomous agents may strain that assumption: they could coordinate
continuously rather than episodically, act in seconds rather than hours, and operate as interacting
populations whose decisions cascade across organizational boundaries. Potential failures include
resource exhaustion from uncoordinated optimization, silent scope violations, cascading pipeline
errors, and patient-facing actions that resist oversight or reversal.

Throughout this Perspective, we follow a hypothetical 68-year-old woman receiving agent-enabled care
for HER2-positive breast cancer. She has diabetes, chronic kidney disease, and treatment-associated
thrombocytopenia. Oncology, symptom-triage, diabetes-management, cardio-oncology, and scheduling agents
each process a portion of her record. In the adverse trajectory used for illustration, fragmented
signals and workflow changes precede a late-detected subdural hematoma. The scenario is not a clinical
prediction or claim of preventability; it is intended to expose coordination questions that can arise
when no actor has responsibility for the joint state of the system.

These are interaction risks: harms that emerge when multiple agents act locally within a shared
clinical environment. A scheduling, oncology, diabetes, and triage agent may each satisfy local safety
criteria while jointly producing unsafe care. Traditional AI evaluation asks whether one model
produces correct outputs under specified conditions; multi-agent deployment additionally asks whether
a population of agents remains safe when actions interact through shared patients, resources, and
workflows. Addressing this question may require infrastructure that can observe agent actions, enforce
boundaries, track dependencies, and intervene before local decisions propagate.

We introduce CIRC, Clinical Infrastructure for Reliable Coordination, a conceptual governance
framework intended to stimulate protocol design rather than a proposed standard. CIRC reframes
clinical agent safety around controls for monitoring, constraining, and coordinating agents as they
scale. This Perspective makes four contributions. First, we identify four architectural gaps that may
emerge when autonomous agents are deployed atop existing interoperability infrastructure. Second, we
define five protocol primitives for monitoring and controlling multi-agent clinical systems. Third, we
propose four descriptive coordination levels, CIRC-L1 through CIRC-L4. Fourth, we argue that
interaction risk should be studied at the population level and outline simulation as one approach for
surfacing it before deployment.

\section{Why existing infrastructure may be insufficient for clinical agents}

Healthcare has invested decades in interoperability infrastructure, including HL7v2 for real-time
clinical messaging, FHIR for structured clinical data exchange, SMART on FHIR for application
authorization, and CDS Hooks for decision support at defined workflow trigger points.\cite{hl7v251,
hl7fhir2019,mandel2016smart,hl7cdshooks2018} These standards are important achievements and could be
adapted to support portions of clinical agent coordination. Missing fields, triggers, or scopes may in
some cases be addressed through profiles, implementation guides, workflow engines, or vendor
extensions. Other gaps concern assumptions about accountable actors, persistent identity, and the
context needed to govern autonomous action across systems.

The broader AI ecosystem has begun addressing agent interoperability through protocols such as
Google's Agent-to-Agent protocol, Anthropic's Model Context Protocol, and IBM's Agent Communication
Protocol.\cite{googlea2a2025,anthropicmcp2024,ibmacp2025} A2A and ACP support communication,
capability discovery, task management, and message exchange, while MCP standardizes how AI systems
connect to tools and data sources. Clinical deployments add constraints that these general protocols
do not themselves standardize. Agents may pursue partially conflicting objectives, and some clinical
failures cannot be safely handled through retry because downstream actions may be difficult to
reverse. The central problem is therefore not only capability discovery but capability constraint:
what an agent may do, for which patients, under which conditions, with which approvals, and how it
must coordinate when its actions intersect with others.

Existing AI risk frameworks characterize what can go wrong. The AI Risk Repository synthesizes more
than 700 risks across 43 taxonomies, and the NIST AI Risk Management Framework organizes governance
functions for trustworthy systems.\cite{slattery2024risk,nist2023airmf} Clinical agent benchmarks
evaluate whether individual agents perform correctly in realistic environments.\cite{jiang2025medagentbench,
ferber2025oncology} Multi-agent clinical deployment additionally raises questions about interacting
agents acting concurrently on shared patients, finite resources, and institutional workflows. Four
gaps illustrate this distinction.

\subsection*{Insufficient agent observability and attribution}

As agents move from recommending actions to initiating orders, updating documentation, routing
referrals, or modifying schedules, audit infrastructure must distinguish human-authored actions,
human-approved agent actions, and autonomous agent-initiated actions. Current mechanisms are commonly
organized around user accounts, applications, service credentials, and record modifications rather
than persistent agent identity across vendors and institutions. FHIR Provenance can represent that a
system, practitioner, organization, or device participated in an activity, but implementations do not
generally require a cross-vendor mechanism for cryptographically verifiable agent identity, model
version, authorization context, and cross-organizational persistence. Without this context,
adverse-event review may struggle to reconstruct which agent acted, under whose authority, and within
which scope.

\subsection*{Missing intent and dependency protocols}

Clinical coordination requires expressing dependencies across disease-focused, patient-focused,
capability-focused, and operational agents. A diabetes-management agent, oncology agent,
record-summarization agent, and imaging-scheduling agent may each impose constraints on one care
pathway. These constraints can include prerequisites, timing windows, resource needs, escalation
rules, data-access limits, and reversal obligations. Healthcare standards can exchange data and invoke
workflow hooks, while general-purpose agent protocols can communicate and hand off tasks. Neither
category, by itself, standardizes a clinical dependency protocol that obliges agents to declare how a
proposed action affects downstream workflows and patient state before execution.

\subsection*{Capability boundaries scoped to human roles}

Healthcare defines role-based authority around expertise, licensure, institutional policy, and
accountability. Agents require analogous constraints, potentially with broad reading privileges but
strict acting privileges. SMART on FHIR provides application-level authorization to data and services,
but these scopes do not encode the complete clinical action context of autonomous behavior. A diabetes
agent might be allowed to recommend insulin adjustments in one setting, require pharmacist approval
in another, and be prohibited from acting autonomously in a third. Clinical deployments therefore may
need enforceable declarations of what an agent can do, for whom, under which conditions, and through
which escalation path.

\subsection*{Limited system-level resource coordination}

When many agents independently conclude that the same time window is optimal for imaging, infusion,
discharge planning, or specialist review, they could collectively overwhelm shared resources even if
each decision is locally reasonable. Prior work calls this type of convergence a ``digital
stampede.''\cite{chopra2025ripple} Historical utilization statistics and single-agent benchmarks may
not expose risks arising dynamically from the interaction of agents, patients, resources, and
workflows. Existing standards can represent many relevant facts, and workflow engines can coordinate
specified processes, but aggregate demand signals and population-level stress tests are not routinely
standardized for clinical agent ecosystems.

\section{The CIRC framework: Levels of Coordination for Clinical Agents}

CIRC can be understood as a set of protocol primitives intended to make agent interactions
observable, bounded, coordinated, and auditable. These primitives are properties of a coordination
layer rather than capabilities delegated to individual agents alone (Table~\ref{tab:primitives}).

\begin{table}[htbp]
\centering
\caption{CIRC protocol primitives in the submitted framework.}
\label{tab:primitives}
\small
\begin{tabularx}{\linewidth}{>{\bfseries}p{0.22\linewidth}YY}
\toprule
Primitive & Definition & Problem addressed \\
\midrule
Persistent agent identity & Clinically meaningful actions carry a verifiable agent identifier and
version tag across organizations. & Attribution, authorization reconstruction, and adverse-event
audit. \\
Structured intent and coordination messages & Machine-readable proposed actions, temporal
constraints, resource needs, clinical prerequisites, and dependencies. & Free-text and
vendor-specific messages resist verification and negotiation. \\
Explicit scope declarations & Signed declarations of authorized clinical scope, data permissions,
and validated populations. & Scope violations may otherwise be visible only after propagation. \\
Dependency tracking & Mechanisms linking actions to downstream obligations across clinical and
administrative systems. & Partial reversal and missing prerequisites can create inconsistent state. \\
Aggregate resource and priority signals & Privacy-preserving summaries of utilization, availability,
and workflow backlog. & Independently optimizing agents can converge on the same constrained
resources. \\
\bottomrule
\end{tabularx}
\end{table}

The submitted framework characterizes CIRC across four levels, with each level adding governance
capabilities and technical sophistication. To ground the levels, it imagines a regional cancer
network of two hospitals and three satellite infusion centers caring for approximately 4,000 patients
receiving systemic therapy over six months. Within this hypothetical network, 180 women receive
adjuvant chemotherapy for HER2-positive breast cancer.\cite{kim2022tchp,hockaday2024tchp} Their care
is orchestrated by oncology-management, diabetes-coordination, cardio-oncology-monitoring, and
scheduling agents.

The recurring patient is a 68-year-old woman with diabetes, chronic kidney disease, and
treatment-associated thrombocytopenia.\cite{kuter2022thrombocytopenia} Over several days, she reports
dizziness and worsening headache. Her platelet count declines across cycles, glucose is elevated, and
wearable data show subtle overnight heart-rate changes. Each signal is processed by a different
agent. The scenario illustrates how infrastructure changes what becomes visible, when review may be
triggered, and how failures might propagate or be contained.

\subsection{CIRC-L1: Isolated agents}

At Level 1, each agent connects to the EHR, scheduling systems, and decision-support tools through
standard interfaces without awareness of concurrent agent activity. Interaction effects are not
directly represented, so resource collisions and cross-domain prerequisite violations may become
visible only downstream.

In the hypothetical trajectory, the triage agent labels dizziness and worsening headache as low
acuity. The oncology agent notes a platelet count of 68,000/$\mu$L and clears chemotherapy without
capturing a decline from 95,000 to 82,000 to 68,000/$\mu$L. The diabetes agent responds to elevated
glucose, while a scheduling agent advances imaging and cancels an observation visit. The example is
constructed to show fragmented state, not to establish that these individual decisions are clinically
correct or that CIRC would have prevented the outcome.

\subsection{CIRC-L2: Vendor-bounded coordination}

At Level 2, agents communicate with agents they know about, typically through common vendors or
explicit integrations. They can exchange structured messages about proposed actions, constraints, and
dependencies, but vendor relationships rather than clinical dependencies determine the coordination
perimeter. This can create coordination islands and asymmetric visibility.

For the hypothetical patient, an oncology agent can object when a same-vendor scheduling agent
proposes moving imaging, but external diabetes, triage, and cardio-oncology agents remain outside the
exchange. The example creates a decision point within one vendor while leaving cross-vendor signals
fragmented.

\begin{figure}[htbp]
\centering
\placeholderfigure{Submitted Figure 1: vendor-bounded coordination simulation}
\caption{CIRC-L2 illustrative conceptual simulation: vendor-bounded coordination. The submitted
caption describes a toy simulation across three vendor ecosystems. It is not based on clinical
deployment, and no conclusions about protocol effectiveness should be drawn.}
\label{fig:l2}
\end{figure}

\subsection{CIRC-L3: Universal protocol-mediated coordination}

At Level 3, a shared coordination protocol enables agents from different vendors and institutions to
exchange structured messages, declare capabilities and constraints, and coordinate over shared
patient state. Governed pathways provide audit trails, constraints, and dependency exchange.

In the hypothetical trajectory, an imaging-schedule adjustment notifies agents with registered
dependencies regardless of vendor. Oncology and diabetes agents raise separate concerns, and a
protocol rule directs the combined concerns to human review. Wearable data can be routed to an agent
or clinician with the relevant scope. These mechanisms create opportunities for review; they do not
establish that the alert is diagnostically valid or clinically effective.

\begin{figure}[htbp]
\centering
\placeholderfigure{Submitted Figure 2: hypothetical error cascade}
\caption{Hypothetical error propagation under different coordination levels. The submitted schematic
contrasts an unchecked upstream error with a protocol-mediated cross-reference checkpoint. It is not
a clinical performance estimate.}
\label{fig:l3}
\end{figure}

\subsection{CIRC-L4: System-aware coordination}

At Level 4, agents become aware of aggregate system state, including utilization, queue dynamics, and
capacity constraints. The coordination layer publishes privacy-preserving summaries without exposing
individual patient information. This could support digital-stampede prevention by allowing proposed
actions to account for constrained resources before rejection.

The submitted framework also associates Level 4 with population pattern detection. In its oncology
example, a system learns that a combination of platelet decline, metabolic instability, reduced
observation, and wearable abnormalities precedes intracranial hemorrhage. Such a function is a
clinical prediction system rather than merely a coordination primitive. It would require independent
development, validation, calibration, subgroup analysis, drift monitoring, and evaluation of false
alerts before influencing individual care.

\begin{figure}[htbp]
\centering
\placeholderfigure{Submitted Figure 3: distributed concern-signal aggregation}
\caption{CIRC-L4 illustrative conceptual simulation: aggregation of distributed concern signals.
The submitted toy example is not calibrated to clinical outcomes and should not be interpreted as an
estimate of sensitivity, specificity, false-positive rate, or prevented harm.}
\label{fig:l4}
\end{figure}

The progression across levels shifts the locus of control from isolated optimization to cross-vendor,
patient-level, and population-level coordination. In the submitted framing, the patient's complexity
is said to demand Level 3 or Level 4 infrastructure. That statement remains conceptual: the framework
does not demonstrate that any level is necessary, sufficient, or superior to alternative
architectures.

\section{Limitations and future directions}

CIRC identifies candidate primitives for multi-agent clinical governance but does not specify a
complete implementation or demonstrate feasibility at scale. Questions remain about cross-vendor
identity, representations for clinically ambiguous intent, adaptation of scope declarations, and
coordination across institutional boundaries. Governance questions include liability, regulatory
classification, institutional autonomy, and local configuration of agent boundaries. These questions
will differ across regimes such as the FDA's device-centered framework and the European Union's
risk-tiered AI Act.\cite{fda2021aiml,eu2024aiact}

The illustrative simulations warrant particular caution. Their parameter choices are not calibrated
against real-world data, populations are small, distributions are predetermined, and outcomes are
known in advance. They illustrate mechanisms by which coordination might interrupt propagation,
expose vendor-boundary gaps, or create review points. They do not estimate clinical sensitivity,
specificity, harm reduction, or real-world performance.

Future work should evaluate CIRC-style controls through progressively stronger evidence. Interaction
risk concerns how agents act jointly on shared patients, resources, and workflows. Population-level
simulation under matched workflow constraints could expose failures missed by single-agent benchmarks,
but it would not establish clinical efficacy. Prospective studies would be needed to determine whether
specific controls improve auditability, reduce unsafe propagation, preserve human oversight, or
prevent system-level failures in clinical settings.

\section{Conclusion}

The hypothetical chemotherapy patient illustrates the central concern of multi-agent clinical AI:
agents can process their local inputs as designed while no infrastructure represents the joint state
of care. The submitted scenario attributes harm to an absence of required sharing, aggregation, and
escalation. This remains an illustrative systems argument rather than evidence that a particular
protocol would have prevented the clinical outcome.

If autonomous agents scale in clinical practice, health systems may need stronger controls for
attribution, structured intent, scope, dependency tracking, and aggregate resource awareness. CIRC
offers one conceptual organization of these concerns. Building and testing coordination mechanisms
before deployment could help institutions surface interaction failures without waiting for patient
harm. The immediate task is to make clinical agent interactions more observable, bounded,
coordinated, and testable while evaluating whether the proposed controls add value beyond existing
standards and alternative architectures.

\section*{Funding}
This work was accomplished without funding.

\section*{Competing interests}
The authors declare no competing interests.

\nocite{*}
\bibliographystyle{unsrtnat}
\bibliography{references}

\end{document}